\begin{lemma}[Claim 4.3(b1): stability of the restoration output]
Let \(\mathcal R\) be a restoration trace satisfying
\noderef{ExtensionBaseRestorationInvariant}, with initial coloring
\(\sigma_n\) and output coloring \(\varrho_1\).  Let \(C\) be a set of
control colors, assume \(v_n\in V(T_j)\), and suppose that \(\sigma_n\) is
\[
 (T(v_n)-v_n,C,\varphi_{j-1})\text{-stable}
\]
in the sense of \noderef{StableColoringBeforeVertex}.  Then
\(\varrho_1\) is
\((T(v_n)-v_n,C,\varphi_{j-1})\)-stable as well.  In the extension-base
application, \(C=D_{j-1}\), and the displayed premise is the current-rung
part (iii) supplied by \noderef{ExtensionRules}.
\end{lemma}
\begin{proof}
Write the trace colorings as
\(\sigma^0=\sigma_n,\sigma^1,\ldots,\sigma^r=\varrho_1\).
Fix a step \(a<r\).  Its invariant says that the switched component
meets \(T_{t_a}\) at exactly \(v_{t_a}\), that \(T_j\) is an ordered
initial segment of \(T_{t_a}\), and that \(v_n\prec v_{t_a}\).
Every vertex in
\(\noderef{TreeVerticesBefore}(T,v_n)\) therefore belongs to
\(T_{t_a}\) and differs from \(v_{t_a}\); hence none of its vertex nodes
belongs to the switched component.

If an edge \(f\) incident with one of those vertices belonged to the
component, its incident vertex node would be adjacent to its edge node
in \noderef{TwoColorIncidenceGraph} and thus would belong to the same
component.  This contradicts the preceding paragraph.  Consequently
\(\sigma^{a+1}(f)=\sigma^a(f)\) for every edge incident with
\(T(v_n)-v_n\).  In particular, all present and missing color sets at
those vertices are unchanged.  Induction over the finite trace gives
\[
 \varrho_1(f)=\sigma_n(f)
 \quad\text{on every such incident edge},\qquad
 \overline{\varrho}_1(w)=\overline{\sigma}_n(w)
 \quad(w\in V(T(v_n)-v_n)).                                  \tag{1}
\]

Now expand the two clauses of
\noderef{StableColoringBeforeVertex}.  The vertex clause for the desired
relation follows by composing (1) with the vertex clause of the assumed
stability of \(\sigma_n\).  If an incident edge has
\(\varphi_{j-1}\)-color in the controlling missing-color union with
\(C\), the edge clause of that same assumed stability first identifies
its \(\sigma_n\)-value with its \(\varphi_{j-1}\)-value, and (1) then
identifies its \(\varrho_1\)-value with the same value.  These are
exactly the two required clauses, proving (b1).
\end{proof}
